\chapter{Persistence and Simulated Disk}

A production consensus system must persist enough state to recover safely
after crashes. For Raft, that includes term, vote, and log state. The Mako
Raft integration includes persistence hooks for Raft log storage and metadata.
This thesis discusses those hooks only as much as needed to explain the
evaluation. It does not claim a full production storage-stack evaluation.

The evaluation uses simulated persistence through the FakeDisk implementation,
which models a queued device with configurable bandwidth and latency. It does
not model every detail of SSD firmware, filesystems, Linux block scheduling, RocksDB
compaction, or device-level writeback. Its purpose is narrower: when
persistence is enabled, it records how much storage work the benchmark
generates.

The instrumentation reports total simulated-disk writes, total simulated-disk bytes,
source-tagged payload bytes, modeled service time, and wait-time diagnostics.
For the thesis body, the most useful metrics are aggregate bytes, write count,
and bytes per committed transaction. Max wait is retained as a diagnostic
counter, but it is a worst-case statistic and is expected to be noisier than
aggregate byte/write counters.

These counters are needed because disk throughput curves alone are not enough.
If Cloud-SSD-like persistence flattens as worker count increases, the thesis
must show whether the benchmark also generated measured persisted data and
write work. Otherwise the result could be dismissed as a benchmark artifact.

\begin{figure}[htbp]
\centering
\includegraphics[width=0.96\linewidth]{figures/manual/fig06_simulated_disk_proof_path.pdf}
\caption[Simulated-disk proof path]{Source-tagged simulated-disk instrumentation. Persistence writes are charged to a simulated disk model and reported as bytes, writes, normalized bytes per transaction, and source breakdowns.}
\label{fig:simulated-disk-proof-path}
\end{figure}
